\section{Frontera de absorción: Interface, Super App y de Chatbot a Agent}

La sección anterior trató los mecanismos internos del Agent; esta sección trata sus fronteras externas. Yao Shunyu (姚顺雨) considera que la mayor oportunidad para las empresas emergentes está en diseñar interfaces distintas. Las capacidades de los modelos pueden dar lugar a formas de interacción beyond ChatGPT y convertirse en una nueva Super App. Aquí, «frontera de absorción» no significa que una empresa vaya a absorberlo todo, sino que cada interfaz de interacción determina en qué tareas puede entrar un sistema inteligente.

\begin{figure}[H]
\centering
\includegraphics[width=0.96\textwidth]{interface-superapp.png}
\caption{Frontera de absorción: la interface determina la oportunidad; distintas formas de interacción pueden dar lugar a distintas Super Apps. Diagrama conceptual propio, elaborado a partir de la conversación entre 00:48:38 y 01:05:01.}
\end{figure}

\begin{knowledgebox}{Lectura de la figura: la interface es una frontera, no solo una UI}
Chatbot, Assistant, Her, interfaces no humanoides, el navegador, el IDE y el sistema operativo pueden convertirse en puntos de entrada a la inteligencia. Cada interface orienta hacia tareas, datos y modelos de negocio distintos, y también moldea la dirección de la investigación.
\end{knowledgebox}

\subsection{Por qué los sistemas Chatbot evolucionan de forma natural hacia Agent}

El sistema Chatbot de una empresa de modelos tiene, por naturaleza, intención del usuario, contexto, necesidad de llamar herramientas y necesidad de estado a largo plazo. Por eso evoluciona de forma natural hacia un sistema Agent: primero recuerda al usuario, luego llama herramientas, después ejecuta tareas y, por último, escribe la retroalimentación de vuelta en el sistema. El Chatbot no es el punto final, sino una de las formas de entrada al Agent.

\begin{figure}[H]
\centering
\includegraphics[width=0.96\textwidth]{chatbot-to-agent.png}
\caption{De Chatbot a Agent: el sistema de chat de una empresa de modelos evoluciona de forma natural hacia un sistema Agent. Diagrama conceptual propio, elaborado a partir de la conversación entre 01:49:28 y 02:05:00.}
\end{figure}

\begin{importantbox}{La transición mínima de Chatbot a Agent}
Basta con que un Chatbot incorpore memory a largo plazo, llamadas a herramientas, estado de la tarea y retroalimentación de resultados para que empiece a convertirse en Agent. La diferencia no está en si la interfaz sigue siendo una ventana de chat, sino en si el sistema puede modificar de manera continua el estado externo.
\end{importantbox}

\subsection{La Super App es un arma de doble filo}

La subsección anterior explicó cómo un Chatbot evoluciona de forma natural hacia un Agent; esta subsección examina la dependencia de trayectoria organizacional que genera un punto de entrada fuerte. Tener una Super App como ChatGPT es una ventaja enorme, porque aporta usuarios, datos, distribución y retroalimentación del producto; pero también moldea la dirección de la investigación y hace que la organización se optimice de forma natural en torno a esa Super App. Si una empresa emergente solo copia una interface existente, difícilmente la superará; si encuentra una interface distinta, puede abrir una nueva frontera de tareas.

\begin{warningbox}{El bloqueo de interfaz influye en la línea de investigación}
Cuando una empresa tiene un punto de entrada fuerte, su investigación tiende a ponerse al servicio de ese punto de entrada. Esto no es negativo, pero la lleva a subestimar las oportunidades de otras interfaces. Históricamente, la sustitución de una plataforma no suele ser una versión mejor de la plataforma anterior, sino una forma de interacción completamente distinta.
\end{warningbox}

\begin{knowledgebox}{Tabla de tipos de Interface}
\begin{tabularx}{\textwidth}{p{0.22\textwidth}p{0.32\textwidth}X}
\toprule
Interface & Tareas típicas & Oportunidad potencial \\
\midrule
Chatbot & Preguntas y respuestas, redacción, explicación, tareas ligeras & Punto de entrada general fuerte, pero propenso a la homogeneización. \\
IDE/Code & Programación, depuración, herramientas de automatización & affordance fuerte; tareas verificables. \\
Browser/OS & Acciones digitales entre páginas web, archivos y aplicaciones & Puede dar lugar a un universal digital agent. \\
Non-human UI & Interacción no humanoide, integrada o en segundo plano & Puede evitar el bloqueo de trayectoria del Chatbot. \\
\bottomrule
\end{tabularx}
\end{knowledgebox}

Por lo tanto, la «frontera de absorción» puede entenderse como la frontera de la interface. Si un modelo solo puede trabajar dentro de una ventana de chat, lo que absorbe son las preguntas y respuestas, la redacción y las tareas ligeras; si entra al IDE, absorbe la producción de código; si entra al navegador y al sistema operativo, absorbe las acciones digitales entre aplicaciones; si entra a los procesos empresariales, absorbe la coordinación repetitiva dentro de la organización. La expansión de la frontera no la logra el modelo por sí solo, sino el modelo, la interfaz, las herramientas y los datos en conjunto.

\subsection{Resumen de la sección}

La frontera de la inteligencia no la determina un solo modelo, sino la interface, las tareas, las herramientas, los hábitos de los usuarios y la trayectoria de la organización en conjunto. Las oportunidades del Agent también surgirán de nuevas formas de interacción.
